\documentclass[11pt]{article}
\usepackage[margin=2.2cm]{geometry}
\usepackage{xcolor}
\usepackage[font=abhaya]{singlish}
\usepackage{booktabs}
\newcommand\demo[1]{\texttt{\detokenize{#1}} & \singlish{#1}}
\title{The \texttt{singlish} package --- sample}
\author{v0.1.0}
\date{}
\begin{document}
\maketitle

\section*{Inline: \texttt{\string\singlish\{...\}}}
\begin{tabular}{ll}
\toprule
Input & Output \\
\midrule
\demo{mama} \\
\demo{amma} \\
\demo{shrii laxkaava} \\
\demo{sixhala} \\
\demo{vidhyaava} \\
\demo{prashnaya} \\
\demo{kSeethraya} \\
\demo{Dhar^maya} \\
\demo{sazda} \\
\demo{nyaShtika bhauthika vidhyaava} \\
\bottomrule
\end{tabular}

\section*{Commands and math pass through}
\singlish{mama \textbf{nyaShtika} \textcolor{blue}{bhauthika vidhyaava} \lat{MESA} paavichchi karanavaa: $E = mc^2$.}

\section*{The \texttt{singlishtext} environment}
\begin{singlishtext}
thaarakaa thula sidhu vana nyaShtika prathikriyaa magin
\textbf{bhaara muuladhravya} nirmaaNaya vee. % (comment stays Latin)
mehi $^{12}\mathrm{C} + {}^{4}\mathrm{He}$ prathikriyaava
vaedhagath vee.
\end{singlishtext}

\section*{Fonts}
\newcommand\fontrow[2]{\texttt{#2} & \singlishsetfont{#1}\singlish{shrii laxkaava -- nyaShtika bhauthika vidhyaava} & \singlishsetfont{#1}\textbf{\singlish{sixhala}}\\}
\begin{tabular}{lll}
\toprule
Option & Regular & Bold \\
\midrule
\fontrow{abhaya}{abhaya}
\fontrow{noto-sans}{noto-sans}
\fontrow{noto-serif}{noto-serif}
\fontrow{yaldevi}{yaldevi}
\fontrow{gemunu}{gemunu}
\bottomrule
\end{tabular}
\end{document}
